\documentclass[11pt,a4paper]{article}

\usepackage[margin=1.05in]{geometry}
\usepackage{amsmath,amssymb,amsthm}
\usepackage{mathtools}
\usepackage{microtype}
\setlength{\emergencystretch}{2.5em}
\usepackage{enumitem}
\usepackage{booktabs}
\usepackage{array}
\usepackage{xcolor}
\usepackage[colorlinks=true,linkcolor=blue!45!black,urlcolor=blue!60!black,citecolor=blue!45!black]{hyperref}

\newtheorem{theorem}{Theorem}[section]
\newtheorem{proposition}[theorem]{Proposition}
\newtheorem{lemma}[theorem]{Lemma}
\newtheorem{corollary}[theorem]{Corollary}
\theoremstyle{definition}
\newtheorem{definition}[theorem]{Definition}
\newtheorem{example}[theorem]{Example}
\theoremstyle{remark}
\newtheorem{remark}[theorem]{Remark}
\newtheorem{warning}[theorem]{Warning}

\newcommand{\RR}{\mathbb{R}}
\newcommand{\CC}{\mathbb{C}}
\newcommand{\NN}{\mathbb{N}}
\newcommand{\ZZ}{\mathbb{Z}}
\newcommand{\norm}[1]{\lvert #1 \rvert}
\newcommand{\dd}{\,d}

\definecolor{okgreen}{rgb}{0.0,0.45,0.15}
\newcommand{\OK}{\textcolor{okgreen}{\textbf{checked}}}

\title{\bfseries $\Gamma$, ${}_0F_1$ and $J_0$\\[0.4em]
\large A short introduction to everything \texttt{list-and-bag.tex} assumes,
starting from $n!$}
\author{}
\date{}

\begin{document}
\maketitle

\noindent
Four things are used without explanation in \texttt{list-and-bag.tex}: the Gamma
function, the reflection formula, the hypergeometric notation ${}_0F_1$, and the
Bessel function $J_0$. This note introduces them in that order, because each one
needs the one before it. Nothing here is original and nothing here is hard; the only
claim the note makes for itself is that it is the shortest path to the four.

\section{$\Gamma$: the factorial, for non-integers}

\subsection{The problem it solves}

We want a function that agrees with $n!$ at the positive integers but is defined --- and
smooth --- in between. Interpolating a sequence is not by itself a well-posed problem:
infinitely many smooth curves pass through the points $1, 1, 2, 6, 24, \dots$. What
picks one out is insisting on the \emph{recursion} rather than just the values.
Factorial satisfies
\[
  (n+1)! = (n+1)\cdot n! ,
\]
so we ask for a function $f$ on $(0,\infty)$ with $f(1) = 1$ and
$f(x+1) = x\,f(x)$ for all $x > 0$.

That is still not enough --- one can perturb a solution by any function of period $1$ ---
but it very nearly is.

\begin{theorem}[Bohr--Mollerup]\label{thm:BM}
There is exactly one function $f : (0,\infty) \to (0,\infty)$ with
\[
  f(1) = 1, \qquad f(x+1) = x f(x), \qquad \log f \text{ convex} .
\]
\end{theorem}

\noindent
Log-convexity is the mild extra demand that rules out the periodic wobbles. The unique
solution is called $\Gamma$, and it is given by an integral.

\subsection{The definition}

\begin{definition}\label{def:gamma}
For $\mathrm{Re}\,z > 0$,
\[
  \Gamma(z) \;=\; \int_0^\infty t^{\,z-1} e^{-t} \dd t .
\]
\end{definition}

Two computations establish everything basic. First, integrating by parts with
$u = t^{z}$ and $dv = -e^{-t}dt$:
\[
  \Gamma(z+1) = \int_0^\infty t^{z}e^{-t}\dd t
  = \underbrace{\bigl[-t^z e^{-t}\bigr]_0^\infty}_{=\,0}
    + z\int_0^\infty t^{z-1}e^{-t}\dd t
  = z\,\Gamma(z) .
\]
Second, $\Gamma(1) = \int_0^\infty e^{-t}dt = 1$. Together these give, by induction,
\begin{equation}\label{eq:factorial}
  \boxed{\;\Gamma(n+1) = n! \quad\text{for } n \in \NN.\;}
\end{equation}
So $\Gamma$ is the factorial with its argument shifted by one. The shift is an
accident of history and is permanently annoying; the thing to memorise is
\eqref{eq:factorial}.

\begin{remark}[why the integral, and not some formula]\label{rem:whyintegral}
Definition~\ref{def:gamma} is not a clever trick for computing $n!$ --- it is
\emph{why $\Gamma$ appears in this subject at all}. Read it right to left: the
integral $\int_0^\infty t^n e^{-t}dt$ \emph{produces} $n!$ out of a power. Whenever a
construction divides by $n!$ and you want to undo that, this integral is the tool,
and that is exactly the role it plays in \texttt{list-and-bag.tex}~\S4: the Borel
transform $\frac{1}{1-x} = \int_0^\infty e^{-t}e^{xt}dt$ is \eqref{eq:factorial}
applied coefficient by coefficient.
\end{remark}

\subsection{The one value that matters here: $\Gamma(\tfrac12) = \sqrt{\pi}$}

\begin{proposition}\label{prop:half}
$\Gamma(\tfrac12) = \displaystyle\int_{\RR} e^{-u^2}\dd u = \sqrt{\pi}$.
\end{proposition}

\begin{proof}
Put $t = u^2$ in Definition~\ref{def:gamma}, so $dt = 2u\,du$ and $t^{-1/2} = u^{-1}$:
\[
  \Gamma(\tfrac12) = \int_0^\infty t^{-1/2}e^{-t}\dd t
  = \int_0^\infty u^{-1} e^{-u^2}\,2u \dd u
  = 2\int_0^\infty e^{-u^2}\dd u
  = \int_{\RR} e^{-u^2}\dd u .
\]
That the last integral is $\sqrt{\pi}$ is the Gaussian integral, proved by squaring it
and passing to polar coordinates.
\end{proof}

This is worth pausing on, because it is the whole hinge of the other note.
\emph{$\Gamma(\tfrac12)$ and the Gaussian integral are the same object.} The
substitution $t = u^2$ is the only thing between them. So when
\texttt{list-and-bag.tex} substitutes $-x^2$ into a generating function and
integrates, it is not doing something to which $\Gamma$ is incidentally relevant: it
is evaluating $\Gamma$ at a half-integer, by construction.

\subsection{Outside the right half-plane}

The integral of Definition~\ref{def:gamma} diverges for $\mathrm{Re}\,z \le 0$, but the
recursion $\Gamma(z) = \Gamma(z+1)/z$ extends $\Gamma$ to all of $\CC$ except
$z = 0, -1, -2, \dots$, where it has simple poles. The useful consequence:

\begin{proposition}\label{prop:entire}
$1/\Gamma$ is entire --- analytic on all of $\CC$, with zeros exactly at
$0, -1, -2, \dots$.
\end{proposition}

This is the fact that licenses the proof of \texttt{list-and-bag.tex}
Theorem~5.1. There the coefficients of $F_s$ are $\varphi(n) = \Gamma(n+1)^{-s}$, and
Ramanujan's master theorem requires $\varphi$ to be analytic at a \emph{negative}
argument, namely $\varphi(-\tfrac12) = \Gamma(\tfrac12)^{-s}$. Had $\Gamma$ itself
been needed there instead of $1/\Gamma$, the pole structure would have obstructed it.

\section{The reflection formula}

\subsection{Statement}

\begin{theorem}[Euler's reflection formula]\label{thm:reflect}
For $z \notin \ZZ$,
\[
  \Gamma(z)\,\Gamma(1-z) \;=\; \frac{\pi}{\sin \pi z} .
\]
\end{theorem}

The name: $z \mapsto 1-z$ reflects the complex plane in the vertical line
$\mathrm{Re}\,z = \tfrac12$, and the formula says $\Gamma$ paired with its own
reflection is elementary. Note the consistency check: the left side has poles at every
integer (from $\Gamma(z)$ at $z \le 0$ and from $\Gamma(1-z)$ at $z \ge 1$), and so
does $\pi/\sin\pi z$. \OK

\subsection{Where it comes from}

A sketch, because the mechanism is the same integral that appears in the other note.
The Beta integral is
\[
  B(a,b) \;=\; \int_0^1 t^{a-1}(1-t)^{b-1}\dd t \;=\; \frac{\Gamma(a)\Gamma(b)}{\Gamma(a+b)} ,
\]
which one proves by writing $\Gamma(a)\Gamma(b)$ as a double integral and changing
variables. Substituting $t = x/(1+x)$ turns it into an integral over $(0,\infty)$:
\[
  \frac{\Gamma(a)\Gamma(b)}{\Gamma(a+b)} \;=\; \int_0^\infty \frac{x^{a-1}}{(1+x)^{a+b}}\dd x .
\]
Now take $b = 1-a$, so that $\Gamma(a+b) = \Gamma(1) = 1$ and the denominator is just
$1+x$:
\begin{equation}\label{eq:keyint}
  \Gamma(a)\Gamma(1-a) \;=\; \int_0^\infty \frac{x^{a-1}}{1+x}\dd x .
\end{equation}
Evaluating \eqref{eq:keyint} by contour integration (a keyhole around the cut on the
positive axis) gives $\pi/\sin \pi a$, which is the theorem.

\subsection{The case we need, and what it explains}

Set $z = \tfrac12$ in Theorem~\ref{thm:reflect}. Then $1 - z = z$, $\sin \tfrac{\pi}{2} = 1$, and
\begin{equation}\label{eq:halfrefl}
  \boxed{\;\Gamma(\tfrac12)^2 \;=\; \pi.\;}
\end{equation}
Which we already knew from Proposition~\ref{prop:half}. The point is not the value but
\emph{which} identity produces it, and that $\tfrac12$ is the fixed point of the
reflection.

\begin{remark}[the $\pi$ in the $\mathrm{List}$ row \emph{is} this formula]\label{rem:tie}
Put $a = \tfrac12$ in \eqref{eq:keyint} and substitute $x = u^2$:
\[
  \Gamma(\tfrac12)\Gamma(\tfrac12) = \int_0^\infty \frac{x^{-1/2}}{1+x}\dd x
  = \int_0^\infty \frac{2\dd u}{1+u^2} = \int_{\RR}\frac{\dd u}{1+u^2} .
\]
So $\int_\RR dx/(1+x^2) = \pi$ is not merely \emph{equal to} $\Gamma(\tfrac12)^2$ ---
it is the reflection formula at $a = \tfrac12$, written in a different variable. This
is the content of \texttt{list-and-bag.tex} Remark~5.2: the exponent $2$ in
$\Gamma(\tfrac12)^2$ comes from reflection pairing $\Gamma(a)$ with $\Gamma(1-a)$, and
the $x^2$ substitution is what sets $a = \tfrac12$, the one place where those two
$\Gamma$-values coincide and the product becomes a \emph{square}. At any other
exponent $x^m$ one gets $\Gamma(\tfrac1m)\Gamma(1-\tfrac1m)$ --- still elementary,
but no longer a power of one number.
\end{remark}

\section{Hypergeometric notation: what ${}_0F_1$ means}

\subsection{Rising factorials}

\begin{definition}\label{def:poch}
The \emph{Pochhammer symbol} (rising factorial) is
\[
  (b)_n \;=\; b(b+1)(b+2)\cdots(b+n-1) \;=\; \frac{\Gamma(b+n)}{\Gamma(b)},
  \qquad (b)_0 = 1 .
\]
\end{definition}

The only instance needed below is $b = 1$:
\begin{equation}\label{eq:one-n}
  (1)_n = 1\cdot 2 \cdots n = n! .
\end{equation}

\subsection{The notation}

Hypergeometric notation is a uniform way of naming power series whose coefficient
\emph{ratios} are rational in $n$. The general shape is
\begin{equation}\label{eq:pFq}
  {}_pF_q\!\left(\begin{matrix} a_1,\dots,a_p \\ b_1,\dots,b_q\end{matrix}; z\right)
  \;=\; \sum_{n \ge 0}
  \frac{(a_1)_n \cdots (a_p)_n}{(b_1)_n\cdots(b_q)_n}\,\frac{z^n}{n!} .
\end{equation}
Read it as: $p$ rising factorials upstairs, $q$ downstairs, and \emph{always} a
spare $n!$ in the denominator. When $p = 0$ the upstairs list is empty and one writes
a bare semicolon in its place: ${}_0F_q(\,;b_1,\dots,b_q;z)$. That is the only feature
of the notation that looks like a typo and is not.

\begin{example}[$p = q = 0$]\label{ex:exp}
${}_0F_0(\,;\,;z) = \sum_n z^n/n! = e^z$. The exponential is the emptiest
hypergeometric series.
\end{example}

\begin{example}[$p=0$, $q=1$, $b_1 = 1$]\label{ex:0F1}
Using \eqref{eq:one-n},
\[
  {}_0F_1(\,;1;z) \;=\; \sum_{n\ge0}\frac{1}{(1)_n}\cdot\frac{z^n}{n!}
  \;=\; \sum_{n \ge 0}\frac{z^n}{n!\cdot n!} \;=\; \sum_{n\ge0}\frac{z^n}{(n!)^2}.
\]
\end{example}

So the single parameter $b_1 = 1$ is doing exactly one job: it contributes a second
$n!$ to the denominator. Repeating it repeats that job, which is the whole reason the
notation shows up in \texttt{list-and-bag.tex}:

\begin{proposition}\label{prop:Fs}
With $F_s(z) = \sum_{n\ge0} z^n/(n!)^s$,
\[
  F_s(z) \;=\; {}_0F_{s-1}\bigl(\,;\underbrace{1,\dots,1}_{s-1};z\bigr) .
\]
In particular $F_1 = {}_0F_0 = e^z$, and $F_2 = {}_0F_1(\,;1;z)$.
\end{proposition}

\begin{proof}
By \eqref{eq:pFq} with $p=0$ and all $b_i = 1$, the $n$-th coefficient is
$1/\bigl((1)_n^{\,s-1}\, n!\bigr) = 1/(n!)^{s-1}n! = 1/(n!)^s$ using \eqref{eq:one-n}.
\end{proof}

\begin{remark}\label{rem:whyF0}
$F_0(z) = \sum z^n = 1/(1-z)$ is \emph{not} of the form ${}_0F_{-1}$ --- the
notation starts at $s = 1$. In hypergeometric terms $F_0$ is
${}_1F_0(1;\,;z)$: one parameter upstairs, $(1)_n = n!$, cancelling the obligatory
$n!$ downstairs. The geometric series sits on the other side of the notation from the
exponential, which is a small notational echo of the $\mathrm{List}$/$\mathrm{Bag}$
distinction the other note is about.
\end{remark}

\section{$J_0$: the Bessel function, and why it is $F_2$}

\subsection{Definition}

\begin{definition}\label{def:J0}
The Bessel function of the first kind of order zero is
\[
  J_0(w) \;=\; \sum_{k \ge 0} \frac{(-1)^k}{(k!)^2}\Bigl(\frac{w}{2}\Bigr)^{2k} .
\]
\end{definition}

It is the solution of $w^2 y'' + w y' + w^2 y = 0$ regular at the origin, and it turns
up whenever a problem has circular symmetry --- the vibration of a drumhead, the
diffraction pattern of a round aperture. For present purposes none of that matters;
only the series does.

\subsection{The identification}

\begin{proposition}\label{prop:J0}
$F_2(-x^2) = J_0(2x)$, and $F_2(z) = I_0(2\sqrt{z})$ where $I_0$ is the modified
Bessel function (Definition~\ref{def:J0} without the $(-1)^k$).
\end{proposition}

\begin{proof}
Substitute $w = 2x$ in Definition~\ref{def:J0}: the $(w/2)^{2k}$ becomes $x^{2k}$, so
\[
  J_0(2x) = \sum_{k\ge0}\frac{(-1)^k x^{2k}}{(k!)^2}
          = \sum_{k\ge0}\frac{(-x^2)^k}{(k!)^2}
          = F_2(-x^2)
\]
by Example~\ref{ex:0F1}. The halving in Definition~\ref{def:J0} is precisely why the
argument is $2x$ and not $x$.
\end{proof}

So the chain of names in \texttt{list-and-bag.tex} unwinds as:
\[
  F_2(z) = \sum \frac{z^n}{(n!)^2}
  \;\xrightarrow[\text{Prop.~\ref{prop:Fs}}]{}\; {}_0F_1(\,;1;z)
  \;\xrightarrow[\text{Prop.~\ref{prop:J0}}]{}\; J_0(2x) \text{ at } z = -x^2 .
\]
Three notations, one series. The reason to mention the third at all is that $J_0$ has
a classical integral that the other two do not advertise:
\begin{equation}\label{eq:J0int}
  \int_0^\infty J_0(t)\dd t = 1
  \qquad\Longrightarrow\qquad
  \int_{\RR} J_0(2x)\dd x = 1 ,
\end{equation}
which is the $s = 2$ case of the theorem, and equals $\Gamma(\tfrac12)^{2-2} = 1$. \OK

\subsection{The convergence caveat, concretely}

For large real $w$,
\[
  J_0(w) \;\sim\; \sqrt{\frac{2}{\pi w}}\,\cos\!\Bigl(w - \frac{\pi}{4}\Bigr) ,
\]
so $J_0$ oscillates with amplitude decaying only like $w^{-1/2}$. That is too slow for
absolute convergence --- $\int \norm{J_0}$ diverges --- but fast enough, with the sign
changes, for \eqref{eq:J0int} to hold conditionally. This is the first point in the
family $F_s$ at which the integral stops being an ordinary convergent integral, and it
is why \texttt{list-and-bag.tex} Remark~5.3 damps by $e^{-\varepsilon x}$ rather than
integrating directly.

Beyond $s = 2$ it gets worse in a specific way. The oscillation of
${}_0F_{s-1}(\,;1,\dots,1;-x^2)$ is not in $x$ but in $x^{2/s}$, so for $s \ge 3$ the
function is not even \emph{periodic} in $x$ --- there is no period to feed to an
oscillatory quadrature rule, and the integral exists only in the regularised
(Mellin/Abel) sense. The $s = 2$ case is the last one where a classical integral
formula like \eqref{eq:J0int} is available, which is why it is the last one verified
numerically in the other note.

\section{Where each of these is used}

\begin{center}
\small
\renewcommand{\arraystretch}{1.25}
\begin{tabular}{@{}llp{0.49\textwidth}@{}}
\toprule
object & here & role in \texttt{list-and-bag.tex} \\
\midrule
$\Gamma(n+1) = n!$ & \eqref{eq:factorial} & the divisor that distinguishes
$\mathrm{Bag}$ from $\mathrm{List}$ \\
$\int_0^\infty t^n e^{-t}dt = n!$ & Rem.~\ref{rem:whyintegral} & the Borel transform
of \S4 --- the de-quotienting map \\
$\Gamma(\tfrac12) = \sqrt{\pi}$ & Prop.~\ref{prop:half} & the Gaussian integral; the
$\mathrm{Bag}$ row \\
$1/\Gamma$ entire & Prop.~\ref{prop:entire} & licenses $\varphi(-\tfrac12)$ in the
proof of Thm.~5.1 \\
$\Gamma(a)\Gamma(1-a) = \pi/\sin\pi a$ & Thm.~\ref{thm:reflect} & the exponent $2$ in
$\Gamma(\tfrac12)^{2-s}$; Rem.~5.2 \\
$a = \tfrac12$ self-dual & \eqref{eq:halfrefl}, Rem.~\ref{rem:tie} & why the
\emph{pair} container $x^2$ and no other power \\
${}_0F_{s-1}$ & Prop.~\ref{prop:Fs} & a name for ``divide by $n!$ exactly $s$ times'' \\
$J_0(2x) = F_2(-x^2)$ & Prop.~\ref{prop:J0} & makes $s=2$ numerically checkable \\
$\int_0^\infty J_0 = 1$ & \eqref{eq:J0int} & the $s=2$ value, $\Gamma(\tfrac12)^0$ \\
$J_0 \sim w^{-1/2}\cos$ & \S4.3 & where convergence stops, and why \\
\bottomrule
\end{tabular}
\end{center}

\section*{Further reading}

\begin{description}[leftmargin=1.6em,style=nextline,nosep]
\item[Artin, \emph{The Gamma Function} (1931).] Thirty-nine pages, and still the best
introduction. Theorem~\ref{thm:BM} is its spine.
\item[Andrews, Askey, Roy, \emph{Special Functions} (1999).] Chapters 1--2 for
$\Gamma$, $B$ and reflection; Chapter 4 for \eqref{eq:pFq}. The standard modern
reference.
\item[Whittaker and Watson, \emph{A Course of Modern Analysis}.] Chapter XII for
$\Gamma$, XVII for Bessel. Where to look when a formula is needed and not understood.
\end{description}

\end{document}
